\section{Results}

\subsection{Feature Coverage}
Table~\ref{tab:features} summarizes the proposed features from the original project goals alongside their implementation status.

\begin{table}[htbp]
\caption{Proposed vs.\ Implemented Features}
\label{tab:features}
\centering
\renewcommand{\arraystretch}{1.3}
\rowcolors{2}{white}{gray!10}
\begin{tabularx}{\columnwidth}{|X|c|}
\hline
\rowcolor[HTML]{4472C4}
\textcolor{white}{\textbf{Feature}} & \textcolor{white}{\textbf{Status}} \\ \hline
CLI accepting \texttt{.ll} files               & \checkmark \\ \hline
LLVM IR parsing via \texttt{llvmlite}          & \checkmark \\ \hline
Typed CFG edge extraction                       & \checkmark \\ \hline
Intermediate scene graph                        & \checkmark \\ \hline
JSON scene graph export                         & \checkmark \\ \hline
Graphviz DOT/PNG export                         & \checkmark \\ \hline
Basic stack animation                           & \checkmark \\ \hline
Rich IR spotlight animation                     & \checkmark \\ \hline
Rich-SSA 3-column animation                     & \checkmark \\ \hline
CFG traversal animation                         & \checkmark \\ \hline
Runtime trace import/export                     & \checkmark \\ \hline
Domtree/loop metadata import                    & \checkmark \\ \hline
GIF output with palette workflow                & \checkmark \\ \hline
Adjustable animation speed                      & \checkmark \\ \hline
Automatic C compilation                         & $\times$   \\ \hline
Interactive stepping/pause                      & $\times$   \\ \hline
YAML configuration files                        & $\times$   \\ \hline
Color palette selection                         & $\times$   \\ \hline
GUI                                             & $\times$   \\ \hline
\end{tabularx}
\end{table}

Of the originally proposed proof-of-concept targets, all were achieved except automatic C-to-IR compilation (users supply \texttt{.ll} files directly) and interactive stepping. The three-column SSA panel and CFG traversal animation were implemented as additional capabilities beyond the original scope.

\subsection{Supported IR Instructions}
LLVManim classifies each LLVM IR instruction into one of nine \texttt{EventKind} categories:

\begin{itemize}
    \item \texttt{alloca}: Stack memory allocation
    \item \texttt{load} / \texttt{store}: Memory read and write operations
    \item \texttt{binop}: Binary arithmetic operations (\texttt{add}, \texttt{sub}, \texttt{mul}, etc.)
    \item \texttt{compare}: Integer and floating-point comparisons (\texttt{icmp},~\texttt{fcmp})
    \item \texttt{call}: Function invocations
    \item \texttt{ret}: Function returns
    \item \texttt{br}: Branch instructions (conditional and unconditional)
    \item \texttt{other}: Unsupported or unrecognized instructions (skipped with a~diagnostic message)
\end{itemize}

CFG edge extraction supports five terminator types: \texttt{br}, \texttt{switch}, \texttt{invoke}, \texttt{indirectbr}, and \texttt{callbr}. Conditional branches are labeled with T/F direction annotations.

\subsection{Animation Modes}
LLVManim provides three stack-focused animation modes and one graph-based mode, each reproducible with a single CLI command. Fig.~\ref{fig:animation_modes} shows representative frames from all three stack modes.

\textbf{Basic mode} (\texttt{--ir-mode basic}) renders a minimal stack-only view. Each IR instruction triggers a per-cell badge flash on the corresponding stack slot, providing a clean visual pulse of execution without additional clutter. This mode is best suited for audiences who need only a high-level sense of execution flow, such as introductory lectures or quick demonstrations, since it strips away all source-level detail and focuses purely on the call stack.
\needspace{2\baselineskip}
\begin{lstlisting}[language=bash,gobble=0]
uv run llvmanim input.ll --animate \
    --ir-mode basic --outdir out/
\end{lstlisting}

\textbf{Rich mode} (\texttt{--ir-mode rich}) extends the basic view with a two-column layout. An IR source panel on the left displays the current function's instructions while a spotlight cursor tracks the executing line, giving the audience both the ``what'' (the stack) and the ``where'' (the source). This mode is particularly valuable when the presenter wants to walk through specific IR constructs---for example, showing how a series of \texttt{alloca} and \texttt{store} instructions initialize local variables before a loop body executes. The spotlight cursor provides a natural focal point that guides the viewer's attention through the code.
\begin{lstlisting}[language=bash,gobble=0]
uv run llvmanim input.ll --animate \
    --ir-mode rich --outdir out/
\end{lstlisting}

\textbf{Rich-SSA mode} (\texttt{--ir-mode rich-ssa}) takes the rich layout a step further with a dedicated third column for SSA values (IR Source $\mid$ SSA Values $\mid$ Stack). The center panel displays the symbolic result of each binary operation, comparison, or load as it executes, making it possible to follow data dependencies without mentally resolving \texttt{\%}-names. This is the most detailed mode and targets audiences already familiar with SSA form---for instance, a compiler course where students need to see how \texttt{\%4 = icmp slt i32 \%2, \%3} evaluates during a loop iteration.
\begin{lstlisting}[language=bash,gobble=0]
uv run llvmanim input.ll --animate \
    --ir-mode rich-ssa --outdir out/
\end{lstlisting}

\begin{figure*}[!htbp]
\centering
\begin{subfigure}[t]{0.49\textwidth}
  \centering
  \includegraphics[width=\linewidth]{rich_mode_frame}
  \caption{Rich mode: two-column layout with IR source spotlight alongside the stack.}
  \label{fig:rich_mode}
\end{subfigure}\hfill
\begin{subfigure}[t]{0.49\textwidth}
  \centering
  \includegraphics[width=\linewidth]{rich_ssa_mode_frame}
  \caption{Rich-SSA mode: three-column layout showing computed SSA values for \texttt{load} and \texttt{icmp} instructions.}
  \label{fig:rich_ssa_mode}
\end{subfigure}

\vspace{0.5em}

\begin{subfigure}[t]{0.70\textwidth}
  \centering
  \includegraphics[width=\linewidth]{basic_mode_frame}
  \caption{Basic mode: stack-only view with per-cell badge flash during a \texttt{call} to \texttt{@init}.}
  \label{fig:basic_mode}
\end{subfigure}
\caption{The three stack-focused animation modes, shown with cropped frames from actual LLVManim output.}
\label{fig:animation_modes}
\end{figure*}

\textbf{CFG traversal} (\texttt{--cfg-animate}) renders an animated walk through the control-flow graph, using Graphviz DOT layout for node positions and a runtime trace for the traversal order:
\begin{lstlisting}[language=bash,gobble=0]
uv run llvmanim input.ll --cfg-animate \
    --dot-cfg main.dot \
    --import-trace trace.json --outdir out/
\end{lstlisting}

\subsection{Test Coverage and Quality}
The LLVManim implementation is validated by 302 automated tests organized into four marker categories:

\begin{table}[htbp]
\caption{Test Distribution by Category}
\label{tab:tests}
\centering
\renewcommand{\arraystretch}{1.3}
\rowcolors{2}{white}{gray!10}
\begin{tabularx}{\columnwidth}{|X|c|}
\hline
\rowcolor[HTML]{4472C4}
\textcolor{white}{\textbf{Category}} & \textcolor{white}{\textbf{Count}} \\ \hline
Unit (ingest, transform)        & 192 \\ \hline
Integration (present, cli, pipeline) & 96 \\ \hline
Contract (export boundaries)    & 8   \\ \hline
End-to-end (entrypoints)        & 6   \\ \hline
\textbf{Total}                  & \textbf{302} \\ \hline
\end{tabularx}
\end{table}

Line coverage stands at 84\% across the entire \texttt{llvmanim} package. Fig.~\ref{fig:coverage_by_module} shows per-module coverage and Fig.~\ref{fig:tests_per_module} shows the test distribution across modules. A four-stage quality gate runs on every pull request:
\begin{enumerate}
    \item \textbf{Ruff}: Linting and import sorting
    \item \textbf{Pyright}: Static type checking (strict mode)
    \item \textbf{import-linter}: Architectural boundary enforcement between layers
    \item \textbf{pytest}: Full test suite execution with coverage reporting
\end{enumerate}

\begin{figure}[htbp]
\centerline{\includegraphics[width=\linewidth]{01_coverage_by_module}}
\caption{Test coverage by module. Most modules exceed the 80\% threshold; \texttt{render/stack\_renderer} (48\%) and \texttt{render/cfg\_renderer} (62\%) are the primary gaps due to Manim rendering dependencies.}
\label{fig:coverage_by_module}
\end{figure}

\begin{figure}[htbp]
\centerline{\includegraphics[width=\linewidth]{02_tests_per_module}}
\caption{Test count per module (298 total). CLI, transform, and ingest modules carry the largest test loads.}
\label{fig:tests_per_module}
\end{figure}

\subsection{Codebase Size}
The LLVManim codebase comprises 3{,}189 lines of Python across 18 modules. Fig.~\ref{fig:lines_of_code} shows the breakdown by layer and module. The Render layer accounts for 42\% of the code, followed by Ingest (25\%), Transform (19\%), and CLI (14\%).

\begin{figure*}[!htbp]
\centerline{\includegraphics[width=\textwidth]{03_lines_of_code}}
\caption{Lines of code by layer (left) and per module (right). The Render layer is the largest due to the Manim animation logic in \texttt{stack\_renderer} and \texttt{cfg\_animation}.}
\label{fig:lines_of_code}
\end{figure*}

\subsection{Performance}
All pipeline stages exhibit linear scaling behavior, confirming that LLVManim can handle realistically sized LLVM IR inputs. Benchmarks use synthetic IR programs with controlled instruction and block counts (median of 5 runs per data point).

Fig.~\ref{fig:pipeline_latency} shows the relative wall-clock cost of each pipeline stage for a 50-block CFG. The Ingest stage (\texttt{parse\_module\_to\_events}) dominates at 11.6\,ms, while all Transform stages complete in under 0.25\,ms. Fig.~\ref{fig:parse_scaling} confirms that parse time scales linearly with instruction count---$\Theta(n)$ where $n$ is the number of IR instructions.

\begin{figure}[htbp]
\centerline{\includegraphics[width=\linewidth]{10_pipeline_stage_latency}}
\caption{Pipeline stage latency for a 50-block CFG (median of 5 runs). Parsing dominates; Transform stages are two orders of magnitude faster.}
\label{fig:pipeline_latency}
\end{figure}

\begin{figure}[htbp]
\centerline{\includegraphics[width=\linewidth]{05_parse_scaling}}
\caption{\texttt{parse\_module\_to\_events} scaling. Linear in instruction count---$\Theta(n)$.}
\label{fig:parse_scaling}
\end{figure}

The scene-graph construction stages also scale linearly: \texttt{build\_scene\_graph(mode='cfg')} is $\Theta(n_b)$ in block count (Fig.~\ref{fig:cfg_build_scaling}), \texttt{derive\_cfg\_trace} is $\Theta(n_b)$ for linear-chain CFGs and $O(n_b \cdot k)$ in the presence of loops (Fig.~\ref{fig:trace_scaling}), and \texttt{build\_scene\_graph(mode='stack')} exhibits sub-linear growth empirically (Fig.~\ref{fig:stack_build_scaling}).

\begin{figure}[htbp]
\centerline{\includegraphics[width=\linewidth]{06_scene_graph_cfg_scaling}}
\caption{CFG scene-graph build time vs.\ block count---$\Theta(n)$.}
\label{fig:cfg_build_scaling}
\end{figure}

\begin{figure}[htbp]
\centerline{\includegraphics[width=\linewidth]{07_trace_scaling}}
\caption{CFG trace derivation scaling---$\Theta(n)$ for linear chains, $O(n \cdot k)$ with loops.}
\label{fig:trace_scaling}
\end{figure}

\begin{figure}[htbp]
\centerline{\includegraphics[width=\linewidth]{08_scene_graph_stack_scaling}}
\caption{Stack scene-graph build time vs.\ call depth---$O(d)$ worst case, sub-linear empirically.}
\label{fig:stack_build_scaling}
\end{figure}

Peak memory allocation also grows linearly with IR size for both pipelines (Fig.~\ref{fig:memory_scaling}), confirming no hidden quadratic blowup in the data structures.

\begin{figure}[htbp]
\centerline{\includegraphics[width=\linewidth]{09_memory_scaling}}
\caption{Peak memory allocation vs.\ CFG block count. Both pipelines scale linearly---$\Theta(n)$.}
\label{fig:memory_scaling}
\end{figure}

\subsection{Reproducibility}
All LLVManim results can be reproduced from the public repository\cite{llvmanim-repo}:

\begin{lstlisting}[language=bash,gobble=0]
git clone https://github.com/BenHoule/LLVManim.git
cd LLVManim
uv sync --dev
./scripts/quality-check.sh  # runs all 4 gates
\end{lstlisting}

System requirements: Python 3.12+, \texttt{uv}, Linux packages \texttt{libcairo2-dev} and \texttt{libpango1.0-dev} (for Manim), and optionally \texttt{graphviz} for PNG graph export.
